% GENERATED FILE. Edit canonical lessons, then run npm run generate:books.
% canonical-source-hash: fnv1a64:9074cdf99b564c81
% canonical-lessons: AR-W47-alif-hamza-below, AR-W47-alif-maqsura, AR-W47-fa, AR-W47-qaf

\chapter{Letters From Goodbyes and Pardons}
\label{ch:ar-letters-from-goodbyes}
\hlchaptermodality{47}

\begin{chapteropening}
\textbf{By the end of this chapter:} \emph{I can write alif with hamza below, alif maqsura, fa and qaf, and find each one in a goodbye or a pardon I already say.}

Complete the last lesson of chapter 47: i can write alif with hamza below, alif maqsura, fa and qaf, and find each one in a goodbye or a pardon I already say.
\end{chapteropening}

\section[alif with hamza below]{\ar{إ} (alif with hamza below) — the alif that opens \ar{إلى}}

\label{lesson:AR-W47-alif-hamza-below}



\begin{quote}
Before the new one: write \ar{ص} and \ar{د} once each, and name them.

Now say \textbf{\ar{إلى} \ar{اللقاء}} (\emph{ilā l-liqāʾ}), \textquotedblleft{}see you\textquotedblright{}. You have read it for a long time. This lesson takes one letter out of it and puts it in your hand.
\end{quote}

\subsection*{Script you'll notice: \ar{إ}}
\textbf{\ar{إ}} — \emph{alif} with a \textbf{hamza below}, read \emph{i} at the start of a word.

It opens \textbf{\ar{إلى}} (\emph{ilā}), \textquotedblleft{}until\textquotedblright{}, in \textbf{\ar{إلى} \ar{اللقاء}} (\emph{ilā l-liqāʾ}), \textquotedblleft{}see you\textquotedblright{}. You already write \textbf{\ar{ا}} and the hamza \textbf{\ar{ء}}. Put the small hamza \textbf{under} the alif and the word opens on \emph{i}; put it over the alif (\textbf{\ar{أ}}) and it opens on \emph{a} or \emph{u}.

\subsection*{Writing: \ar{إ}}
Put your pen on \ar{إ} and follow its line. Copy the shape you can see — slowly, and larger than it is printed.

\begin{quote}
This book does not yet tell you \textbf{where to start this character or which way to travel}. It is not written down here until it can be written down with a source. Copying what is in front of you needs no such source, and it is how the shape gets into your hand in the meantime.
\end{quote}

\subsection*{Your turn}
\begin{itemize}
\raggedright
  \item \emph{Look:} at \textbf{\ar{إلى} \ar{اللقاء}} and point at \ar{إ} inside it
  \item \emph{Trace it:} \ar{إ} three times, saying \emph{i} as you finish each one
  \item \emph{Write it:} \ar{د}, then \ar{إ}, from memory
  \item \emph{Read it:} \textbf{\ar{إلى} \ar{اللقاء}} aloud once more
\end{itemize}

\begin{tcolorbox}[breakable,colback=teal!4,colframe=teal!35!black,title={Before you move on}]
Which letter is this — \ar{إ}? (\textbf{alif with hamza below}, \emph{i}.) Which word did it come from? (\textbf{\ar{إلى} \ar{اللقاء}}, \emph{ilā l-liqāʾ}, \textquotedblleft{}see you\textquotedblright{}.)
\end{tcolorbox}
\section[alif maqṣūra]{\ar{ى} (alif maqṣūra) — the letter that ends \ar{إلى}}

\label{lesson:AR-W47-alif-maqsura}



\begin{quote}
Before the new one: write \ar{د} and \ar{إ} once each, and name them.

Now say \textbf{\ar{إلى} \ar{اللقاء}} (\emph{ilā l-liqāʾ}), \textquotedblleft{}see you\textquotedblright{}. You have read it for a long time. This lesson takes one letter out of it and puts it in your hand.
\end{quote}

\subsection*{Script you'll notice: \ar{ى}}
\textbf{\ar{ى}} — \emph{alif maqṣūra}, a long \emph{ā} at the very end of a word.

It ends \textbf{\ar{إلى}} (\emph{ilā}) in \textbf{\ar{إلى} \ar{اللقاء}}. It looks like \textbf{\ar{ي}} with its two dots taken away, and that is the whole trick: no dots, no \emph{y}. It only ever stands at the end of a word, and there it is read as a long \emph{ā}.

\subsection*{Writing: \ar{ى}}
\hlblockfigure{\detokenize{figures/AR-W47-alif-maqsura-filmstrip.pdf}}{How \ar{ى} is written, stroke by stroke}

\begin{itemize}
\raggedright
  \item \textbf{1.} start here: curve from above the baseline through the upper half of the S
  \item \textbf{2.} without lifting, continue through the wide flat lower curve and finish near the baseline — and only now lift
\end{itemize}

\textbf{Pen lifts: 0.} The pen never leaves the paper.

\begin{quote}
Stroke order is one attested teaching order; hands and teachers vary. Source: Arabic Language Learning Notes, Learn How to Write Arabic Letters — Basic Naskh: Alif Maqsura \ar{ى}, isolated form and stroke order (updated 2025-01-29; accessed 2026-08-23).
\end{quote}

\subsection*{Your turn}
\begin{itemize}
\raggedright
  \item \emph{Look:} at \textbf{\ar{إلى} \ar{اللقاء}} and point at \ar{ى} inside it
  \item \emph{Trace it:} \ar{ى} three times, saying \emph{ā} as you finish each one
  \item \emph{Write it:} \ar{إ}, then \ar{ى}, from memory
  \item \emph{Read it:} \textbf{\ar{إلى} \ar{اللقاء}} aloud once more
\end{itemize}

\begin{tcolorbox}[breakable,colback=teal!4,colframe=teal!35!black,title={Before you move on}]
Which letter is this — \ar{ى}? (\textbf{alif maqṣūra}, \emph{ā}.) Which word did it come from? (\textbf{\ar{إلى} \ar{اللقاء}}, \emph{ilā l-liqāʾ}, \textquotedblleft{}see you\textquotedblright{}.)
\end{tcolorbox}
\section[fāʾ]{\ar{ف} (fāʾ) — the letter inside \ar{عفوا}}

\label{lesson:AR-W47-fa}



\begin{quote}
Before the new one: write \ar{إ} and \ar{ى} once each, and name them.

Now say \textbf{\ar{عفوا}} (\emph{ʿafwan}), \textquotedblleft{}you're welcome\textquotedblright{}. You have read it for a long time. This lesson takes one letter out of it and puts it in your hand.
\end{quote}

\subsection*{Script you'll notice: \ar{ف}}
\textbf{\ar{ف}} — \emph{fāʾ}, the sound \emph{f}.

It is the second letter of \textbf{\ar{عفوا}} (\emph{ʿafwan}), \textquotedblleft{}you're welcome\textquotedblright{}. The letter is a small closed \textbf{head} on a broad bowl, with \textbf{one dot above}.

\subsection*{Writing: \ar{ف}}
\hlblockfigure{\detokenize{figures/AR-W47-fa-filmstrip.pdf}}{How \ar{ف} is written, stroke by stroke}

\begin{itemize}
\raggedright
  \item \textbf{1.} start here: loop counterclockwise around the small closed head
  \item \textbf{2.} continue left through the broad bowl without lifting
  \item \textbf{3.} lift once, then place the upper dot last — and only now lift
\end{itemize}

\textbf{Pen lifts: 1.}

\begin{quote}
Stroke order is one attested teaching order; hands and teachers vary. Source: Introduction to Arabic: Egyptian Arabic for first-year students, Alphabet \ar{ف} \ar{ق}, independent Faa at 00:01.7–00:03.3 (University of Oregon Libraries Open Textbook, 2023; accessed 2026-08-23).
\end{quote}

\subsection*{Your turn}
\begin{itemize}
\raggedright
  \item \emph{Look:} at \textbf{\ar{عفوا}} and point at \ar{ف} inside it
  \item \emph{Trace it:} \ar{ف} three times, saying \emph{f} as you finish each one
  \item \emph{Write it:} \ar{ى}, then \ar{ف}, from memory
  \item \emph{Read it:} \textbf{\ar{عفوا}} aloud once more
\end{itemize}

\begin{tcolorbox}[breakable,colback=teal!4,colframe=teal!35!black,title={Before you move on}]
Which letter is this — \ar{ف}? (\textbf{fāʾ}, \emph{f}.) Which word did it come from? (\textbf{\ar{عفوا}}, \emph{ʿafwan}, \textquotedblleft{}you're welcome\textquotedblright{}.)
\end{tcolorbox}
\section[qāf]{\ar{ق} (qāf) — the letter inside \ar{اللقاء}}

\label{lesson:AR-W47-qaf}



\begin{quote}
Before the new one: write \ar{ى} and \ar{ف} once each, and name them.

Now say \textbf{\ar{إلى} \ar{اللقاء}} (\emph{ilā l-liqāʾ}), \textquotedblleft{}see you\textquotedblright{}. You have read it for a long time. This lesson takes one letter out of it and puts it in your hand.
\end{quote}

\subsection*{Script you'll notice: \ar{ق}}
\textbf{\ar{ق}} — \emph{qāf}, a \emph{k} made far back in the throat.

It is inside \textbf{\ar{اللقاء}} (\emph{al-liqāʾ}), \textquotedblleft{}the meeting\textquotedblright{}, in \textbf{\ar{إلى} \ar{اللقاء}}. Its head is the same as \textbf{\ar{ف}}, which you wrote in the last lesson, but it carries \textbf{two dots} and, standing alone, its bowl dips deeper, below the line.

\subsection*{Writing: \ar{ق}}
\hlblockfigure{\detokenize{figures/AR-W47-qaf-filmstrip.pdf}}{How \ar{ق} is written, stroke by stroke}

\begin{itemize}
\raggedright
  \item \textbf{1.} start here: loop counterclockwise around the small closed head
  \item \textbf{2.} continue down and left through the deep bowl without lifting
  \item \textbf{3.} lift once, then place the upper-right dot
  \item \textbf{4.} lift again, then place the upper-left dot — and only now lift
\end{itemize}

\textbf{Pen lifts: 2.}

\begin{quote}
Stroke order is one attested teaching order; hands and teachers vary. Source: Introduction to Arabic: Egyptian Arabic for first-year students, Alphabet \ar{ف} \ar{ق}, independent Qaf at 00:01.5–00:03.5 (University of Oregon Libraries Open Textbook, 2023; accessed 2026-08-23).
\end{quote}

\subsection*{Your turn}
\begin{itemize}
\raggedright
  \item \emph{Look:} at \textbf{\ar{إلى} \ar{اللقاء}} and point at \ar{ق} inside it
  \item \emph{Trace it:} \ar{ق} three times, saying \emph{q} as you finish each one
  \item \emph{Write it:} \ar{ف}, then \ar{ق}, from memory
  \item \emph{Read it:} \textbf{\ar{إلى} \ar{اللقاء}} aloud once more
\end{itemize}

\begin{tcolorbox}[breakable,colback=teal!4,colframe=teal!35!black,title={Before you move on}]
Which letter is this — \ar{ق}? (\textbf{qāf}, \emph{q}.) Which word did it come from? (\textbf{\ar{إلى} \ar{اللقاء}}, \emph{ilā l-liqāʾ}, \textquotedblleft{}see you\textquotedblright{}.)
\end{tcolorbox}
